% Chemistry Homework Formatter — student paper preamble
% Default wrapper is English (CIE papers). Load ctex only if the teacher asks for Chinese.
% Compiler: XeLaTeX. Overleaf: upload this file as main.tex plus the images/ folder.
\documentclass[10pt,a4paper]{article}

\usepackage[margin=16mm,headheight=13pt,headsep=5pt,footskip=9mm]{geometry}
\usepackage{fontspec}
\usepackage[version=4]{mhchem}
\usepackage{chemfig}
\usepackage{graphicx}
\usepackage{xcolor}
\usepackage{booktabs}
\usepackage{array}
\usepackage{enumitem}
\usepackage{needspace}
\usepackage{setspace}
\usepackage{fancyhdr}
\usepackage{amsmath,amssymb}

\raggedbottom
\setlist[enumerate]{leftmargin=1.6em,itemsep=0.25em,topsep=0.25em}
\setlist[itemize]{leftmargin=1.4em,itemsep=0.12em}

\pagestyle{fancy}
\fancyhf{}
\fancyhead[L]{\small \homeworktitle}
\fancyhead[R]{\small \homeworkmarks}
\fancyfoot[C]{\thepage}
\renewcommand{\headrulewidth}{0.3pt}

\newcommand{\homeworktitle}{Chemistry homework}
\newcommand{\homeworkmarks}{}
\newcommand{\sethomework}[2]{%
  \renewcommand{\homeworktitle}{#1}%
  \renewcommand{\homeworkmarks}{#2}%
}

\newcommand{\answerlines}[1]{%
  \par\vspace{0.3em}%
  \foreach \n in {1,...,#1}{\noindent\rule{\linewidth}{0.2pt}\\[1.05em]}%
}

% Fallback if pgffor is missing: comment \answerlines and use \vspace instead.
\usepackage{pgffor}

\graphicspath{{images/}}
\DeclareGraphicsExtensions{.png,.jpg,.jpeg,.pdf}

% Measure the image, then break only if the whole labelled clip will not fit.
\newcommand{\clipq}[3]{%
  \par
  \sbox0{\includegraphics[width=0.94\linewidth,keepaspectratio]{#3}}%
  \Needspace{\dimexpr\ht0+1.35em+2pt\relax}%
  \noindent{\small\textbf{#1}}\hfill{\small[#2]}\par\vspace{0.08em}%
  \noindent\usebox0\par
  \vspace{0.28em}%
}

\newcommand{\qfigure}[1]{%
  \sbox0{\includegraphics[width=0.94\linewidth,keepaspectratio]{#1}}%
  \Needspace{\dimexpr\ht0+0.6em\relax}%
  \noindent\usebox0\par
}

\newcommand{\qfigurebig}[1]{%
  \sbox0{\includegraphics[width=\linewidth,keepaspectratio]{#1}}%
  \Needspace{\dimexpr\ht0+0.6em\relax}%
  \noindent\usebox0\par
}

\setlength{\parindent}{0pt}
\setlength{\parskip}{0.12em}

\setlist[enumerate]{leftmargin=1.55em,itemsep=0.06em,topsep=0.12em}
\newcommand{\qhead}[1]{%
  \par\vspace{0.38em}%
  {\normalsize\bfseries #1}\par\vspace{0.12em}%
}
\newcommand{\mcq}[1]{%
  \par\noindent\textbf{#1.}\enspace\ignorespaces
}
\newcommand{\mcqend}{\hfill{[1]}\par}
\newcommand{\optrow}[4]{%
\par\noindent{\small
\begin{tabular}{@{}p{0.48\linewidth}@{}p{0.48\linewidth}@{}}
\textbf{A}\; #1 & \textbf{B}\; #2 \\[0.12em]
\textbf{C}\; #3 & \textbf{D}\; #4
\end{tabular}}\par
}
\newcommand{\optlist}[4]{%
\par
\begin{enumerate}[label=\textbf{\Alph*.}, leftmargin=1.55em, itemsep=0.05em, topsep=0.08em]
\item #1
\item #2
\item #3
\item #4
\end{enumerate}
}
\newcommand{\qfig}[2][0.58]{%
\par
\begin{center}\vspace{-0.1em}%
\includegraphics[width=#1\linewidth,keepaspectratio]{#2}%
\end{center}\vspace{-0.1em}%
}

\begin{document}
\sethomework{AS stoichiometry homework 3 · empirical and molecular formulae}{18 marks}

\begin{center}
{\large\bfseries AS stoichiometry homework 3 · empirical and molecular formulae}\\[0.2em]
{\small CIE 9701 AS Chemistry · 18 MCQ · 18 marks}
\end{center}

\noindent Name: \underline{\hspace{5.8cm}} \quad Class: \underline{\hspace{2.6cm}} \quad Date: \underline{\hspace{2.8cm}}

\vspace{0.25em}
{\small Topics: 8 empirical formulae · 9 molecular formulae · 10 hydrates · 11 combustion analysis · 12 formulae and equations. Later questions also use earlier skills.

Circle \textbf{one} letter (A--D) for each question. Show working. Calculators and a Periodic Table may be used.}
\vspace{0.15em}

\noindent\begin{minipage}{\linewidth}
\qhead{8 · Empirical formulae}
\mcq{1}%
What is the empirical formula of butanoic acid?
\optrow{\ce{C2H4O}}{\ce{C3H6O}}{\ce{C4H8O}}{\ce{C5H10O}}
\mcqend
\end{minipage}\vspace{0.16em}

\noindent\begin{minipage}{\linewidth}
\mcq{2}%
Compound X consists of 40.0\% carbon, 6.7\% hydrogen and 53.3\% oxygen by mass.

What is the empirical formula of compound X?
\optrow{\ce{CH2O}}{\ce{C2H2O}}{\ce{C2H4O}}{\ce{CHO}}
\mcqend
\end{minipage}\vspace{0.16em}

\noindent\begin{minipage}{\linewidth}
\mcq{3}%
Compound X is an organic compound that contains 30.6\% carbon, 3.8\% hydrogen, 20.4\% oxygen and 45.2\% chlorine by mass.

What is the empirical formula of X?
\optrow{\ce{C2H3OCl}}{\ce{C2H4OCl}}{\ce{C3H4OCl}}{\ce{C4H3O2Cl2}}
\mcqend
\end{minipage}\vspace{0.16em}

\noindent\begin{minipage}{\linewidth}
\mcq{4}%
R is an oxide of Period 3 element T. 5.00\,g of R contains 2.50\,g of T.

What is T?
\optlist{magnesium}{aluminium}{silicon}{sulfur}
\mcqend
\end{minipage}\vspace{0.16em}

\noindent\begin{minipage}{\linewidth}
\qhead{9 · Molecular formulae}
\mcq{5}%
Compound X is a straight chain hydrocarbon with an $M_\mathrm{r}$ of 84.

What can be determined about X?
\begin{enumerate}[label=\textbf{\arabic*.}, leftmargin=1.6em]
\item empirical formula
\item molecular formula
\item whether X contains a \ce{C=C} bond or not
\end{enumerate}
The responses \textbf{A} to \textbf{D} should be selected on the basis of

{\centering
{\small
\begin{tabular}{cccc}
\toprule
\textbf{A} & \textbf{B} & \textbf{C} & \textbf{D} \\
\midrule
1, 2 and 3 are correct & 1 and 2 only are correct & 2 and 3 only are correct & 1 only is correct \\
\bottomrule
\end{tabular}}\par}

{\small No other combination of statements is used as a correct response.}
\mcqend
\end{minipage}\vspace{0.16em}

\noindent\begin{minipage}{\linewidth}
\mcq{6}%
An aqueous solution contains 4.00\,g of a carboxylic acid, Q. When this solution reacts with an excess of magnesium, 380\,cm$^3$ of gas is produced, measured at s.t.p.

What is the relative formula mass of Q?
\optrow{59}{118}{126}{236}
\mcqend
\end{minipage}\vspace{0.16em}

\noindent\begin{minipage}{\linewidth}
\mcq{7}%
P is a compound that burns in an excess of oxygen to give carbon dioxide and water only.

2.20\,g of P contains 1.20\,g of carbon and 0.20\,g of hydrogen.

When P is added to a solution of sodium carbonate, bubbles of gas are seen.

What is P?
\optlist{\ce{CH3CHO}}{\ce{CH3CO2H}}{\ce{CH3COCH2CH2OH}}{\ce{CH3CH2CH2CO2H}}
\mcqend
\end{minipage}\vspace{0.16em}

\noindent\begin{minipage}{\linewidth}
\qhead{10 · Hydrated and anhydrous compounds}
\mcq{8}%
A sample of 35.6\,g of hydrated sodium carbonate contains 25.84\% sodium ions by mass.

When this sample is heated, anhydrous sodium carbonate and water are formed.

Which mass of water is given off?
\optrow{7.2\,g}{10.6\,g}{14.4\,g}{21.2\,g}
\mcqend
\end{minipage}\vspace{0.16em}

\noindent\begin{minipage}{\linewidth}
\mcq{9}%
Hydrated cobalt(II) sulfate loses water when heated to give anhydrous cobalt(II) sulfate. All the water of crystallisation is lost to the atmosphere as steam.

When 3.10\,g of hydrated cobalt(II) sulfate, \ce{CoSO4.xH2O}, is heated to constant mass the \textbf{loss} in mass is 1.39\,g.

What is the value of $x$, to the nearest whole number?
\optrow{4}{6}{7}{11}
\mcqend
\end{minipage}\vspace{0.16em}

\noindent\begin{minipage}{\linewidth}
\mcq{10}%
An experiment is carried out to determine the value of $x$ in hydrated lithium hydroxide, \ce{LiOH.xH2O}. A sample of the solid is heated in a crucible over a Bunsen flame.

Complete dehydration takes place; decomposition does \textbf{not} occur.
\begin{align*}
\text{mass of empty crucible\,/\,g} &= Q\\
\text{mass of crucible with \ce{LiOH.xH2O}\,/\,g} &= R\\
\text{mass of crucible and residue after heating\,/\,g} &= S
\end{align*}
Which equation gives the correct value of $x$?
\optlist{$\dfrac{S-R}{18}$}{$\dfrac{R-S}{18}$}{$\dfrac{23.9(R-S)}{18(S-Q)}$}{$\dfrac{23.9(S-R)}{18(S-Q)}$}
\mcqend
\end{minipage}\vspace{0.16em}

\noindent\begin{minipage}{\linewidth}
\qhead{11 · Combustion analysis}
\mcq{11}%
Compound X contains the elements C, H and O only.

2.00\,g of X produces 4.00\,g of carbon dioxide and 1.63\,g of water when completely combusted.

What is the empirical formula of X?
\optrow{\ce{CHO2}}{\ce{C2H2O}}{\ce{C2H4O}}{\ce{CH2O2}}
\mcqend
\end{minipage}\vspace{0.16em}

\noindent\begin{minipage}{\linewidth}
\mcq{12}%
Substance Q is a hydrocarbon. When 1.00\,g of Q is completely burned, 3.22\,g of carbon dioxide is produced.

What could be the identity of Q?
\optlist{cyclohexene}{cyclopentane}{ethene}{pentane}
\mcqend
\end{minipage}\vspace{0.16em}

\noindent\begin{minipage}{\linewidth}
\mcq{13}%
When a small sample of hydrocarbon Q is completely combusted, it produces 3.52\,g of carbon dioxide and 1.44\,g of water.

What could be the structure of hydrocarbon Q?
\qfig[0.78]{fig-q13-structures.png}
\mcqend
\end{minipage}\vspace{0.16em}

\noindent\begin{minipage}{\linewidth}
\qhead{12 · Chemical formulae and equations}
\mcq{14}%
Which pair of formulae is correct?
\optlist{\ce{Ag2CO3} and \ce{(NH4)3NO3}}{\ce{K2HCO3} and \ce{Zn3(PO4)2}}{\ce{AgHCO3} and \ce{K3PO4}}{\ce{ZnCO3} and \ce{(NH4)2PO4}}
\mcqend
\end{minipage}\vspace{0.16em}

\noindent\begin{minipage}{\linewidth}
\mcq{15}%
Chlorine dioxide, \ce{ClO2}, reacts with aqueous sodium hydroxide to produce water and a mixture of two sodium salts, \ce{NaClO2} and \ce{NaClO3}.

What is the mole ratio of \ce{NaClO2} to \ce{NaClO3} in the product mixture?
\optrow{1:2}{3:5}{1:1}{5:3}
\mcqend
\end{minipage}\vspace{0.16em}

\noindent\begin{minipage}{\linewidth}
\qhead{Earlier skills · combining masses, reacting masses, equations}
\mcq{16}%
Originally, chemists thought indium oxide had the formula \ce{InO}. By experiment they showed that 4.8\,g of indium combined with 1.0\,g of oxygen to produce 5.8\,g of indium oxide. The $A_\mathrm{r}$ of oxygen was known to be 16.

Which value for the $A_\mathrm{r}$ of indium is calculated using these data?
\optrow{38}{77}{115}{154}
\mcqend
\end{minipage}\vspace{0.16em}

\noindent\begin{minipage}{\linewidth}
\mcq{17}%
A 4.00\,g sample of an anhydrous Group 2 metal nitrate, Z, is heated strongly until there is no further change of mass. A solid residue of mass 1.37\,g is formed.

Which metal is present in Z?
\optlist{barium}{calcium}{magnesium}{strontium}
\mcqend
\end{minipage}\vspace{0.16em}

\noindent\begin{minipage}{\linewidth}
\mcq{18}%
Copper dissolves in dilute nitric acid producing a blue solution of \ce{Cu(NO3)2}, water and nitrogen(II) oxide as the only products.

How many moles of acid react with three moles of copper in the balanced equation?
\optrow{2}{4}{6}{8}
\mcqend
\end{minipage}\vspace{0.16em}

\end{document}
